% !TEX program = lualatex
% acta.tex: acta de un pleno. Ejemplo con nombres inventados.
\documentclass[11pt,a4paper]{article}
\usepackage{acta-consejo}
% config.tex: datos fijos de los documentos (órgano, firmante y marca BORRADOR).
% Es lo único que hay que cambiar al empezar un curso nuevo.
\setorgano{el}{Consejo de Estudiantes}
\renewcommand{\firmante}{Nombre Apellidos}
\renewcommand{\cargofirmante}{Secretario \delorgano{} de la Facultad de Ciencias}
% Mientras el documento no sea definitivo, deja \borradortrue: añade la marca
% «BORRADOR» de fondo. Para la versión que se firma, \borradorfalse.
\borradortrue


\begin{document}

\begin{paracol}{2}
% ------------------------------------------ Columna izquierda: asistencia
\begin{asistencia}
\grupo{Comisión Permanente}
\asis[Grado en Física]{Lucía Ejemplo Pérez}
\asis[Grado en Química]{Mario Modelo Gil}
\asis[Grado en Matemáticas]{Nombre Apellidos}

\grupo{Grado en Biotecnología}
\asis{Ana Prueba Sanz}
\asis{Jorge Muestra Lázaro}

\grupo{Grado en Geología}
\asis{Clara Ficticia Royo}

\grupo{Excusan su asistencia}
\asis{Pablo Inventado Ruiz}
\end{asistencia}

\switchcolumn
% ------------------------------------------ Columna derecha: apertura
En Zaragoza, siendo las 17:30 del jueves 1 de enero de 2099, en el Salón de
Actos del Edificio B de la Facultad de Ciencias, en segunda convocatoria, se
reúne el Pleno \delorgano{} con la asistencia de las personas que figuran al
margen y el siguiente orden del día:

\punto{Aprobación, si procede, del acta anterior.}
Se aprueba por asentimiento.

\punto{Informe de la Comisión Permanente.}
Dña. Lucía Ejemplo resume las actividades realizadas desde el último pleno.

\end{paracol}

% ------------------------------------------ Resto del acta a ancho completo
\punto{Elecciones de representantes.}
\subpunto{Comisiones de Garantía de la Calidad.}
Se presentan las candidaturas y se acuerda votarlas de forma telemática.

\subpunto{Delegaciones de titulación.}
Se elige por asentimiento a Dña. Ana Prueba como delegada de Biotecnología.

\punto{Ruegos y preguntas.}
No hay.

\cierre{19:05}
\firma

\end{document}
